\section{Global Stokes Lean Anchor Alignment}
\label{sec:global_stokes_lean_anchor_alignment}
\begingroup
\sloppy
\emergencystretch=3em

This fragment records the current Lean anchors that refine the open nodes in
\Cref{sec:global_assembly}.  Each entry is meant to be usable as a small
coordination task: instantiate the named Lean package, discharge the listed
explicit hypotheses, then replace the corresponding planned blueprint node by
the implemented anchor.

\begin{definition}[Interior selected and extended chart boxes]
  \label{anc:interior_chart_box_anchors}
  \lean{Stokes.interiorChartSelectedBox, Stokes.interiorChartExtendedBox,
    Stokes.interiorChartExtendedBox.mk_of_subsets}
  \leanok
  \uses{def:chart_derivative_transition, thm:chartwise_smooth_transport}
  \emph{Carries blueprint nodes:}
  \texttt{def:interior\_chart\_selected\_box} and
  \texttt{def:interior\_chart\_extended\_box}.
  The selected-box anchor records \(a \le b\), the closed coordinate box inside
  the chart-transition domain, and support of
  \texttt{ManifoldForm.transitionPullbackInChart} inside that box.  The
  extended-box anchor adds an open neighborhood on which the same chart
  representative is \texttt{ContDiffOn}.

  \emph{Remaining explicit hypotheses:}
  concrete lower and upper corners; target and overlap inclusions; support
  containment for the chart representative being localized; and the
  neighborhood smoothness statement.  The local Euclidean Stokes theorem for
  these interior boxes is still a separate task.
\end{definition}

\begin{definition}[Selected box partition of unity]
  \label{anc:selected_box_partition_of_unity}
  \lean{Stokes.SelectedBoxPartitionOfUnity,
    Stokes.SelectedBoxPartitionOfUnity.mkOfBoxes,
    Stokes.SelectedBoxPartitionOfUnity.mkOfSubsets,
    Stokes.exists_smoothPartitionOfUnity_subordinate_chartAt_source}
  \leanok
  \uses{anc:interior_chart_box_anchors}
  \emph{Carries blueprint node:}
  \texttt{def:partition\_subordinate\_selected\_boxes}.  It also supplies the
  finite activity bridge needed by \texttt{lem:partition\_piece\_support}
  through \texttt{mem\_active\_of\_mem\_fintsupport} and
  \texttt{mem\_active\_of\_mem\_tsupport}.

  \emph{Remaining explicit hypotheses:}
  a smooth partition \(\rho\); a compact control set \(K\); a finite active set;
  the proof that every partition term active on \(K\) lies in that finite set;
  selected extended boxes for all active indices; and, for existence, the
  compactness argument selecting boxes from chart sources.  The current package
  is interior-chart shaped; adding boundary-box indices or a unified
  interior/boundary sum wrapper remains open.
\end{definition}

\begin{definition}[Partition-localized form data]
  \label{anc:localized_form_data}
  \lean{Stokes.ManifoldForm.localizedForm, Stokes.LocalizedFormData,
    Stokes.LocalizedFormData.mkLocalized,
    Stokes.LocalizedFormData.interiorChartSelectedBox}
  \leanok
  \uses{anc:selected_box_partition_of_unity, def:inchart_transition_pullback}
  \emph{Carries blueprint nodes:}
  \texttt{def:partition\_localized\_form},
  \texttt{lem:partition\_piece\_support}, and the interior part of
  \texttt{def:chart\_supported\_form}.  The anchor stores the scalar
  coefficient, the localized form, the equality with pointwise scalar
  multiplication, and the selected-box support data needed to pass a localized
  form to interior chart-box APIs.

  \emph{Remaining explicit hypotheses:}
  smoothness of the scalar partition coefficient; the support proof for
  \(\rho_i \omega\) in the selected box; compatibility with boundary selected
  boxes; the finite-sum reconstruction of \(\omega\); and the exterior
  derivative finite-sum and Leibniz rules
  \texttt{lem:extderiv\_finite\_sum} and
  \texttt{lem:partition\_piece\_extderiv}.
\end{definition}

\begin{definition}[Boundary orientation map bridge]
  \label{anc:boundary_chart_orientation_map_data}
  \lean{Stokes.BoundaryChartOrientationMapData,
    Stokes.BoundaryChartOrientationMapDataOn,
    Stokes.boundaryChartOrientationCompatibleOn_of_orientationMapDataOn,
    Stokes.boundaryChartPreservesOrientationOn_of_orientationMapDataOn}
  \leanok
  \uses{def:boundary_chart_oriented_atlas, thm:boundary_orientation_from_atlas}
  \emph{Carries blueprint nodes:}
  the orientation-map content behind
  \texttt{def:boundary\_chart\_oriented\_atlas} and
  \texttt{thm:boundary\_orientation\_from\_atlas}; it is also the clean bridge
  into \texttt{thm:boundary\_transition\_open\_image},
  \texttt{thm:boundary\_change\_of\_variables}, and
  \texttt{thm:boundary\_chart\_integral\_invariance}.

  \emph{Remaining explicit hypotheses:}
  a tangential linear equivalence for each boundary-coordinate point; equality
  of its underlying linear map with
  \texttt{boundaryChartTransitionTangentMap}; either a positive determinant or
  an \texttt{Orientation.map} equality; and setwise data upgrading the
  pointwise bridge to the selected boundary face.  Local inverse data and image
  equality for boundary boxes are still supplied by the boundary-chart
  image-data packages.
\end{definition}

\begin{definition}[Project-local integral wrappers]
  \label{anc:project_local_integral_wrappers}
  \lean{Stokes.projectLocalBulkIntegral, Stokes.projectLocalBoundaryIntegral,
    Stokes.projectLocalStokes_of_boundaryChartExtendedBox,
    Stokes.projectLocalStokes_of_orientedAtlas_imageData,
    Stokes.projectLocalStokes_of_orientedManifold_imageData}
  \leanok
  \uses{thm:boundary_chart_local_stokes, thm:boundary_chart_integral_invariance}
  \emph{Carries blueprint nodes:}
  \texttt{def:project\_local\_bulk\_integral},
  \texttt{def:project\_local\_boundary\_integral}, and the boundary-chart part
  of \texttt{thm:project\_local\_chart\_supported\_stokes}.  The wrappers fix
  the project-facing names for the local bulk and outward-first boundary terms.

  \emph{Remaining explicit hypotheses:}
  selected or extended boundary-chart boxes; oriented-atlas or oriented-manifold
  assumptions when transporting the boundary term to a target chart; and
  boundary image data identifying the source lower face with the target
  boundary box.  Interior project-local wrappers have not yet been added.
\end{definition}

\begin{definition}[Finite-sum global assembly data]
  \label{anc:global_stokes_assembly_data}
  \lean{Stokes.GlobalStokesAssemblyData, Stokes.globalStokesAssembly,
    Stokes.sum_localChartSupportedStokes}
  \leanok
  \uses{anc:project_local_integral_wrappers}
  \emph{Carries blueprint nodes:}
  \texttt{thm:sum\_local\_stokes},
  \texttt{thm:global\_bulk\_eq\_local\_sum}, and
  \texttt{thm:global\_boundary\_eq\_local\_sum}.  This is the algebraic
  finite-cover layer: local Stokes identities for interior and boundary pieces
  are summed, artificial interior boundary terms are cancelled, and the result
  is compared with the stored global bulk and boundary real numbers.

  \emph{Remaining explicit hypotheses:}
  the finite active chart set; interior and boundary piece finsets; bulk and
  boundary real-valued term functions; local Stokes identities for every active
  piece; reconstruction of the global bulk and boundary integrals from those
  sums; and cancellation of artificial interior faces.  The field
  \texttt{chartChangeInvariance} is presently only recorded as a proposition;
  the newer final package below makes the chart-change sum equality explicit.
\end{definition}

\begin{definition}[Final global Stokes data package]
  \label{anc:global_stokes_data}
  \lean{Stokes.GlobalStokesData, Stokes.GlobalStokesData.stokes,
    Stokes.globalStokes}
  \leanok
  \uses{anc:global_stokes_assembly_data,
    anc:boundary_chart_orientation_map_data}
  \emph{Carries blueprint node:}
  \texttt{thm:global\_compact\_support\_stokes} as the current final Lean
  theorem anchor.  It refines the older assembly package by naming a separate
  boundary partition sum and by making chart-change compatibility an equality
  between boundary-chart terms and boundary-partition terms.

  \emph{Remaining explicit hypotheses:}
  all fields of \texttt{GlobalStokesData}: finite chart and piece data; local
  Stokes identities for interior and boundary pieces; global bulk
  reconstruction; artificial interior-boundary cancellation; chart-change
  cancellation into boundary partition terms; and global boundary
  reconstruction.  Compact support, manifold-with-boundary assumptions,
  partition-of-unity construction, and actual manifold integral definitions are
  not in the theorem statement yet; they must be used to build an instance of
  this data package.
\end{definition}

\begin{definition}[Project-local final Stokes package]
  \label{anc:project_local_global_stokes}
  \lean{Stokes.ProjectLocalGlobalStokesData,
    Stokes.ProjectLocalGlobalStokesData.stokes,
    Stokes.projectLocalGlobalStokes}
  \leanok
  \uses{anc:project_local_integral_wrappers, anc:global_stokes_data}
  \emph{Carries blueprint nodes:}
  \texttt{thm:project\_local\_chart\_supported\_stokes},
  \texttt{thm:sum\_local\_stokes}, and the project-local specialization of
  \texttt{thm:global\_compact\_support\_stokes}.  This is the most direct
  finite-sum target for boundary-chart boxes because the local terms are
  exactly \texttt{projectLocalBulkIntegral} and
  \texttt{projectLocalBoundaryIntegral}.

  \emph{Remaining explicit hypotheses:}
  active chart and piece finsets; source and target chart choices for every
  piece; lower and upper box corners; the equality reducing the represented
  global bulk integral to project-local bulk sums; pointwise project-local
  Stokes for every active piece; chart-change cancellation into boundary
  partition terms; and boundary partition reconstruction.  A useful next
  collaboration task is to instantiate \texttt{localProjectStokes} from
  \texttt{projectLocalStokes\_of\_boundaryChartExtendedBox} or the oriented
  image-data variants.
\end{definition}

\begin{theorem}[Compact and mathlib-facing final statements]
  \label{anc:remaining_final_global_nodes}
  \lean{Stokes.globalStokes}
  \leanok
  \uses{anc:selected_box_partition_of_unity, anc:localized_form_data,
    anc:project_local_global_stokes}
  \emph{Carries blueprint nodes conditionally:}
  \texttt{thm:compact\_manifold\_stokes} and
  \texttt{thm:compare\_mathlib\_integral\_api} should not be marked closed by
  the bare algebraic theorem.  Instead, they should point to
  \texttt{Stokes.globalStokes} as the final assembly target once the data
  package is produced from compactness, partitions of unity, boundary
  orientation, and the chosen integral API.

  \emph{Remaining explicit hypotheses:}
  conversion from compact manifold hypotheses to compact support data; a
  canonical or project-local global integral definition; comparison with any
  future mathlib manifold-form integral; and a proof that the selected-box
  partition/localized-form layer supplies every field of either
  \texttt{GlobalStokesData} or \texttt{ProjectLocalGlobalStokesData}.
\end{theorem}

\endgroup
